\section{Introduction}
\label{sec:introduction}

Large language models (LLMs) have become a common component in modern
recommendation systems. Recent methods use LLM-derived embeddings,
sequential adapters, graph representations, profile modules, or intent
representations to improve candidate scoring~\cite{llm2rec,llmemb,llmesr,irllrec,proex,promax,elmrec,rlmrec}.
These systems are normally judged by ranking metrics, but comparisons across
them are difficult to trust: different papers use different candidate sets,
different backbones, different importers, and different evaluation schemas, so
it is rarely clear whether a reported gain reflects the modeling idea or the
evaluation setup. A prerequisite for studying any LLM-based scoring signal is
therefore a comparison in which the candidate set, backbone, and evaluation
pipeline are held fixed.

We address this with a rigorous unified same-candidate reranking protocol. We
focus on the reranking stage: each method receives the same user histories and
the same 101 candidates per event (one ground-truth positive and 100
popularity-sampled negatives), and all outputs are converted into the same
score schema before metrics are computed. Every method that requires an LLM
uses a single shared Qwen3-8B backbone, a shared importer, exact coverage
checks, and paired Holm-corrected bootstrap tests, with visible provenance.
This same-candidate protocol intentionally avoids claims about full-catalog
retrieval; it isolates the question of which scoring signal ranks a fixed
candidate set most accurately. We instantiate the protocol over eight Amazon
domains (Beauty, Books, Electronics, Movies, Sports, Toys, Home, Tools),
10{,}000 users each (Beauty 973), against the same eight official-code-level
LLM-based recommendation baselines in every domain.

Within this protocol, our central empirical finding is simple. A
\emph{task-grounded pointwise LLM relevance posterior} -- the model's own
per-candidate relevance probability, queried directly and used as the ranking
score -- ranks first in \emph{six of the eight} domains (Books, Electronics,
Sports, Toys, Home, Tools), improving NDCG@10 by $+21.6\%$ to $+53.2\%$ over the
strongest baseline; it leads on all seven metrics (HR@5/@10/@20,
NDCG@5/@10/@20, MRR) in five of these six domains (Electronics, Sports, Toys,
Home, Tools; 56/56 paired-significant) and on six of seven on Books (HR@20
versus LLMEmb the exception). It is not first in the other two:
on Beauty it is rank~2, $-11\%$ on NDCG@10 behind ProEx, though its shortfall
versus ProEx is not Holm-significant on any metric, leaving it statistically
indistinguishable from the strongest baseline; on Movies it is mid-pack,
rank~5 and $-24\%$ behind LLMEmb, substantially behind the four stronger
baselines. We state these two losses openly rather than
restricting the reported domain set to the wins. A random floor and a popularity
baseline both sit near NDCG@10~$\approx 0.045$, so the gains are over a
genuinely informative comparison, not the floor. This says that for
same-candidate reranking, a directly elicited pointwise posterior from a strong
general-purpose LLM is a remarkably strong signal in most domains, ahead of
trained-embedding and adapter-based recommenders run under identical conditions,
while leaving room for stronger baselines where the protocol favours
profile- or embedding-based structure (Beauty, Movies).

This work originally set out to study \emph{actionable uncertainty}: whether a
calibrated candidate posterior, instrumented with a task-grounded uncertainty
decomposition (boundary ambiguity, a calibration gap, evidence support, and
counterevidence) and a validation-controlled risk-adjusted ranking family
($\eta=0$ recovering the bare posterior), could rank a fixed candidate set more
reliably. We therefore report a rigorous \emph{negative result} about this
design space. Under a leave-one-component-out ablation and a raw-probability
attribution (run on the four domains with full per-event signal rows: Sports,
Toys, Home, Tools), the uncertainty decomposition and risk penalty do not
improve same-candidate ranking: the zero-risk setting ($\eta=0$) is test-best in
all four diagnostic domains, a confidence-only variant (the calibrated posterior
with no uncertainty and no risk) matches or exceeds the full family in three of
those four domains, and the raw probability without any calibration map tracks
the full family within $\leq 0.008$ (mean $\approx 0.005$) NDCG@10 on the three
domains with local raw-probability rows (Sports, Home, Tools); the fourth
diagnostic domain's row is server-side. The boundary-uncertainty term is
inert in all four, and removing the counterevidence and risk terms is
non-worse. The working ranking signal is located in the pointwise posterior
itself, not in the uncertainty machinery built on top of it. As a minor,
descriptive observation, the event-level uncertainty signal does stratify
ranking reliability across all four diagnostic domains, but this is partly
circular and we do not claim it improves ranking. We accordingly present the
calibrated-uncertainty/risk family as an analyzed and characterized design
space, not as a validated headline contribution.

Our contributions are:

\begin{enumerate}
    \item \textbf{A unified same-candidate reranking protocol and benchmark.}
    We build an eight-domain, 10{,}000-user (Beauty 973), 101-candidate (one
    positive, 100 popularity-sampled negatives) same-candidate benchmark with
    eight official-code-level adapted LLM recommendation baselines, a single
    shared Qwen3-8B backbone, a shared importer, exact coverage checks, full
    HR/NDCG/MRR reporting, paired Holm-corrected bootstrap tests, and visible
    provenance (Table~\ref{tab:baseline_provenance}). This lets LLM-based
    scoring signals be compared without candidate-set, backbone, or
    importer confounds.

    \item \textbf{A task-grounded pointwise LLM posterior ranks first in six of
    eight domains.} Under this protocol, the model's directly elicited
    per-candidate relevance probability ranks first in six of the eight domains,
    improving NDCG@10 by $+21.6\%$ to $+53.2\%$ over the strongest baseline and
    leading on all seven metrics in five of those six domains (Electronics,
    Sports, Toys, Home, Tools; 56/56 paired-significant) and on six of seven on
    Books (HR@20 versus LLMEmb the exception). We report the
    two non-winning domains explicitly (Beauty rank~2, $-11\%$; Movies rank~5,
    $-24\%$) instead of dropping them. We compute paired Holm-corrected bootstrap
    tests on all eight domains: on the six winning domains \method{} is
    significantly ahead on all 56 metric/baseline pairs (54/56 on Books), while on
    the two losing domains the paired tests quantify the gap honestly --- on Beauty
    none of the losses to ProEx is Holm-significant, and on Movies \method{} is
    significantly behind the stronger baselines on 22 of 56 pairs. The claim is
    scoped to same-candidate reranking, not
    full-catalog SOTA.

    \item \textbf{A rigorous negative result on calibrated uncertainty and
    risk.} Via leave-one-component-out ablations and a raw-probability
    attribution, we show that the calibrated-uncertainty decomposition and the
    risk-adjusted ranking family do not improve same-candidate ranking
    ($\eta=0$ test-best in all domains; confidence-only $\geq$ full family in
    three of four; components inert or mildly harmful), locating the working
    signal in the posterior. We keep the uncertainty signal's descriptive
    reliability stratification as a caveated minor observation, not a ranking
    claim.
\end{enumerate}
